% !PRINTED 81
\weberpage{81}
%<<CONT

von \(a_\nu^{(\mu)}a_k^{(i)}\) in der Entwickelung von \(A\) und \(A_\nu^{(\mu)}\) der Coëfficient von \(a_\nu^{(\mu)}\), folglich \(A_{\nu,k}^{\mu,i}\) der Coëfficient von \(a_k^{(i)}\) in der Determinante \(A_\nu^{(\mu)}\).

Man kann nach diesem Satze die sogenannte geränderte Determinante
\begin{equation}
        U=\begin{vmatrix}
        a_1^{(1)}&a_2^{(1)}&\cdots&a_n^{(1)}&u_1\\
        a_1^{(2)}&a_2^{(2)}&\cdots&a_n^{(2)}&u_2\\
        \vdots&\vdots&&\vdots&\vdots\\
        a_1^{(n)}&a_2^{(n)}&\cdots&a_n^{(n)}&u_n\\
        v_1&v_2&\cdots&v_n&q
        \end{vmatrix}
\tag{11}
\end{equation}
nach den Elementen der letzten Zeile und Colonne entwickeln und erhält:
\begin{equation}
        U=qA-\sum^{i}\sum^{k} u_i v_k A_k^{(i)}.
\tag{12}
\end{equation}

Man erhält diese Gleichung aus (10), wenn man \(n\) in \(n+1\) verwandelt, und die Elemente der letzten Zeile und Colonne durch eine andere Bezeichnung auszeichnet.

Auch bei den höheren Unterdeterminanten ist bisweilen die Bezeichnung durch Differentialquotienten zweckmässig, so dass z.~B.
\begin{equation}
        A_{\nu,k}^{i,\mu}=
        {\partial^2 A\over \partial a_\nu^{(i)}\,\partial a_k^{(\mu)}}
\tag{13}
\end{equation}
gesetzt wird.




\begin{center}
{\large \S.\ 24.}

\vspace{0.5em}
{\large\gesperrt{Lineare homogene Gleichungen.}}
\end{center}

Die hauptsächlichste Anwendung der Determinanten, der die ganze Theorie ihren Ursprung verdankt, ist die Auflösung linearer Gleichungen.

Wir wollen hier die Aufgabe gleich in allgemeinster Weise in Angriff nehmen, da die specielle Form kaum eine Vereinfachung ist und sich nachher leicht aus dem allgemeinen Resultate ableiten lässt.

Wir betrachten ein System von $m$ Gleichungen ersten Grades, in denen $n$ Unbekannte $x_1,x_2,\ldots,x_n$ homogen vorkommen:
%CONT>>
